\begin{abstract}
Se clasifican 104 categorías de flores de la competencia \textit{Petals to the
Metal} con redes entrenadas desde cero a $112\times112$, usando 12,753 imágenes
de entrenamiento y 3,712 de validación, con un desbalance de 18 a 782 imágenes
por clase. A partir de una CNN de referencia de cuatro convoluciones, se cambia
un factor a la vez: una MLP como línea base, una CNN con dos convoluciones más,
la eliminación de la regularización, pesos por clase en la pérdida, una tasa de
aprendizaje menor y menos dropout. La CNN de referencia es la mejor, con F1 macro
de 0.5784 y exactitud de 64.06\% en validación. La MLP no logra aprender (F1 de
0.018) y la red más profunda no mejora (0.566). Sin regularización, la brecha de
F1 entre entrenamiento y validación sube de 0.235 a 0.344, y los pesos por clase
bajan el F1 a 0.488 porque sacrifican las clases frecuentes por las raras. Los
resultados indican que la limitación principal es la cantidad de datos por clase,
no la capacidad de la red. Cada configuración usa una sola semilla; el código,
los historiales y un notebook que reproduce los resultados están en un
repositorio público.
\end{abstract}

\begin{keywords}
Clasificación de imágenes \sep Redes neuronales convolucionales \sep
Regularización \sep Desbalance de clases \sep F1 macro \sep Flores
\end{keywords}
